\section{Introduction}
\label{sec:introduction}

Measuring productive capacity is fundamental to macroeconomic theory, yet capacity output is unobservable. Conventional output-gap and survey measures impose strong cyclical or technical assumptions, often conflating engineering limits with demand fluctuations. \citet{Shaikh2016} proposes an influential alternative: an accounting-based measure derived from a long-run cointegrating regression between aggregate output and capital stock in the US corporate sector (1947--2011). Shaikh interprets the resulting long-run coefficient as an empirical anchor for potential output and treats the residuals as cyclical capacity utilization. This paper critically replicates Shaikh's empirical procedure to evaluate whether the estimated output--capital relation represents a stable technical benchmark or a specification-sensitive artifact.

To evaluate this relation, I reframe the aggregate output--capital coefficient as a capacity transformation elasticity ($\theta \equiv \partial \ln Y^p / \partial \ln K$). Under the balanced-growth knife-edge assumption standard in post-Keynesian models, this elasticity is constrained to unity ($\theta = 1.0$), meaning productive capacity expands in lockstep with capital accumulation. When $\theta \neq 1.0$, however, capacity formation decouples from accumulation: an elasticity below unity ($\theta < 1.0$) defines structural overaccumulation, where capital accumulates faster than capacity expands, while $\theta > 1.0$ generates chronic excess capacity. Reconstructing Shaikh's single-equation ARDL(2,4) baseline (Stage S0) demonstrates that the estimation is approximately reproducible within canonical data bounds ($\hat{\theta} \approx 0.72$ versus Shaikh's published 0.66), yielding a point estimate within the structural overaccumulation interval.

Yet this baseline point estimate is fragile across model choices. In Stage S1, I evaluate a combinatorial grid of 500 ARDL specifications varying lag lengths, deterministic trends, and structural break subsets. While 102 specifications satisfy the baseline $F$-bounds test at the 10\,\% level (and 62 at the 5\,\% level), the recovered transformation elasticity varies widely across the informational frontier ($\hat{\theta} \in [0.65, 0.95]$). Rather than converging to a unique technical constant, the single-equation framework exhibits disciplined non-uniqueness: the estimated capacity ceiling depends substantially on researcher choices regarding lag profiles and deterministic terms.

Moving from single-equation regressions to a joint system framework reveals a deeper econometric limitation. In Stage S2, I evaluate 48 attempted bivariate Vector Error Correction Models (VECMs, of which 36 are successfully estimated) sweeping lag orders and deterministic configurations over the 1947--2011 sample. The standalone bivariate output--capital relation fails to cointegrate in every tested specification. Residuals exhibit persistent non-stationary drift. Aggregate output and capital stock do not form an empirically self-sufficient cointegrating system in isolation over the post-war period.

Long-run system stability becomes recoverable only when the rate of exploitation enters the state vector alongside historical shock controls. In a trivariate VECM $(\ln Y, \ln K, \ln e)$, where $\ln e_t = \ln \pi_t - \ln \omega_t$ captures functional income distribution, six specifications achieve cointegration rank $r=1$ and companion matrix stability. In the focal model, the estimated elasticity point estimate ($\hat{\theta} = 0.727$) sits in the overaccumulation range, but carries substantial normalized estimation uncertainty ($\text{SE} = 4.852, t = -0.15$). In sharp contrast, the coefficient on the rate of exploitation is identified with high precision ($\hat{\beta}_e = 20.022, \text{SE} = 2.536, t = 7.90$). The system demonstrates that functional distribution enters the long-run attractor with dominant statistical significance, while the transformation elasticity remains imprecisely estimated once unconstrained by single-equation assumptions.

These empirical findings motivate a conceptual reinterpretation of capacity formation. The output--capital relation does not fail because capital is unproductive; rather, productive capacity is distributionally conditioned. In capitalist economies, converting capital equipment into potential output depends on the intensity of the labor process, work organization, and shift arrangements. These parameters are not neutral engineering constraints. They are structured by the distribution of income between profits and wages and the institutional settlements that govern the workplace. System stability emerges only when the econometric model incorporates the distributional mechanisms through which capital accumulation is socially coordinated and disciplined.

The remainder of the paper proceeds as follows. Section~\ref{sec:literature_review} traces the historical shift in capacity utilization metrics from engineering surveys to cyclical disciplining tools. Section~\ref{sec:conceptual_framework} develops the conceptual framework, formalizing the transformation elasticity and its connection to Marxian profit rate decompositions and choice-of-technique dynamics. Section~\ref{sec:replication_strategy} presents the three-stage econometric results: single-equation reconstruction (S0), specification geometry across the ARDL lattice (S1), and system-level VECM estimation (S2). Section~\ref{sec:discussion} discusses the theoretical implications of structural overaccumulation, the institutional interpretation of distributional conditioning, and methodological boundaries. Section~\ref{sec:conclusion} concludes.
